\section{Galois groups and partial local correspondences}\label{sec:preliminaries}

We collect the anabelian inputs and the local calculation that will connect Frobenius elements with cyclotomic characters. The substantial inputs from Uchida and Hoshi are stated explicitly; the remaining arguments use standard global and local class field theory, Kummer theory, and Chebotarev's theorem.

\subsection{Conventions and solvably closed fields}
Fix an algebraic closure $\overline{\Q}$ of $\Q$. For a number field $M\subset\overline{\Q}$ put $G_M=\Gal(\overline{\Q}/M)$, and put
\[
 \mathcal G=G_{\Q},\qquad
 P=\Gal(\Q^s/\Q),
\]
where $\Q^s\subset\overline{\Q}$ is the maximal prosolvable extension of $\Q$. For a number field $M\subset\Q^s$ write $P_M=\Gal(\Q^s/M)$. All homomorphisms between profinite groups and all characters in this paper are continuous. An open homomorphism is a homomorphism with open image: a continuous surjection between profinite groups is a quotient map and is open.

We use finite primes unless infinity is explicitly mentioned. For a finite prime $v$, its absolute norm is $Nv$, and its absolute residue degree is $[\kappa(v):\F_p]$, where $p$ is its rational residue characteristic. Frobenius is arithmetic. A chosen prolongation of $v$ determines a decomposition group $D_v$ and inertia group $I_v$; a different prolongation conjugates both. Expressions involving several such groups always use specified compatible prolongations.

\begin{lemma}\label{lem:solvable-closure}
Every solvably closed algebraic extension $E$ of $\Q$ contains $\Q^s$. The field $\Q^s$ is solvably closed. For every number field $M\subset\Q^s$, the extension $\Q^s/M$ is the maximal prosolvable extension of $M$.
\end{lemma}
\begin{proof}
Since $E(\mu_n)/E$ is abelian, $E$ contains every root of unity. Adjoining an $n$th root of an element of $E$ then gives a cyclic extension, so $E$ is closed under radicals. Finite solvable Galois extensions can be built by towers of radical extensions after adjoining roots of unity; hence $\Q^s\subset E$.

The same tower argument shows that $\Q^s$ contains all roots of unity and is closed under radicals. Any nontrivial finite abelian extension would have a cyclic subextension of prime degree, which Kummer theory excludes. Any nontrivial abelian algebraic extension would have a nontrivial finite abelian subextension. Thus $\Q^s$ is solvably closed.

Finally, a finite solvable Galois extension of $M$ lies in a finite solvable Galois extension of $\Q$: first take a solvable normal closure of $M/\Q$ inside $\Q^s$, and then the compositum of the finitely many conjugate extensions. The group of this compositum over that normal closure embeds in a product of solvable groups. This proves the last assertion.
\end{proof}

In particular, every additive $\ell$-adic character of $G_M$ factors through $P_M$ when $M\subset\Q^s$. For any solvably closed Galois extension $E/F$, the action on roots of unity defines its cyclotomic character
\[
 \chi_{E/F}:\Gal(E/F)\longrightarrow\widehat{\Z}^{\times}.
\]
We use $\chi_{M,\ell}$ for the usual $\ell$-adic character of $G_M$, and $\chi_\ell$ for that of $P$. The $\ell$-adic logarithm is normalized to vanish on the finite torsion subgroup of $\Z_\ell^\times$.

\subsection{The anabelian inputs}
The first input permits a homomorphism to be extended from an open subgroup to the normalizer of its kernel.

\begin{theorem}[Uchida]\label{thm:uchida-isomorphism}
Let $E_i/F_i$ be solvably closed Galois extensions of number fields. Every isomorphism
$\theta:\Gal(E_1/F_1)\xrightarrow{\sim}\Gal(E_2/F_2)$
is induced by a unique field isomorphism $j:E_2\xrightarrow{\sim}E_1$, with
$g j=j\theta(g)$ for all $g$.
\end{theorem}
This is \cite[Theorem]{uchida1979}. One consequence we shall use twice is the following.
\begin{lemma}\label{lem:centralizer}
The centralizer in $P$ of any open subgroup of $P$ is trivial.
\end{lemma}
\begin{proof}
An open subgroup is $P_M$ for a number field $M\subset\Q^s$. An element centralizing $P_M$ is a field automorphism of $\Q^s$ implementing the identity automorphism of $P_M$. The identity field map also implements it, so uniqueness in Theorem~\ref{thm:uchida-isomorphism} gives the assertion.
\end{proof}

Here is the local input, in the precise setting in which we need it. For a homomorphism $\varphi:G_M\to P_N$, an odd prime $\ell$, and a prime $v\nmid\ell$ of $M$, say that $v$ is \emph{active at $\ell$} if $\varphi$ is nontrivial on a pro-$\ell$ Sylow subgroup of $I_v$. The condition is independent of the chosen Sylow subgroup and prolongation, since such choices are conjugate.

\begin{theorem}[Uchida's partial local correspondence]\label{thm:local-correspondence}
Let $M$ be a number field, let $N\subset\Q^s$ be a number field, and let $\varphi:G_M\to P_N$ be open. Fix an odd prime $\ell$. Every prime $v\nmid\ell$ active at $\ell$ determines a unique prime $w\nmid\ell$ of $N$ for which, with suitable prolongations,
\[
 \varphi(D_v)\subset D_w.
\]
Moreover, $\varphi$ is injective on a pro-$\ell$ Sylow subgroup of $I_v$, and sends this subgroup into $I_w$. All but finitely many primes of $N$ arise from such active primes of $M$.
\end{theorem}
This is the odd-prime case of the trichotomy in \cite[\S1, pp.~595--596]{uchida1981}, together with its Proposition~1. In that trichotomy, nontrivial pro-$\ell$ inertia is case~(ii); case~(i) kills that inertia, and the torsion case~(iii) requires $\ell=2$. The open-image hypothesis, rather than surjectivity onto $P_N$, suffices. We will apply the theorem after restricting to smaller open subgroups as well. The cofinite target coverage in this theorem is for a fixed $\ell$; its exceptional set is not asserted to be uniform in $\ell$.

The uniqueness assertion concerns primes of the finite target field $N$. When restricting a map to smaller open subgroups, we will intersect a containing target decomposition group with $P_{N'}$ for a finite extension $N'/N$ inside $\Q^s$, and then use uniqueness at the $N'$-level. Proposition~\ref{prop:frobenius-generation} carries out this comparison with explicit source and target fields; no uniqueness of infinite prolongations is required.

The final input converts cyclotomic compatibility into the required field embedding.
\begin{theorem}[Hoshi]\label{thm:hoshi}
Let $E_i/F_i$ be solvably closed Galois extensions of number fields, and let
$\alpha:\Gal(E_1/F_1)\to\Gal(E_2/F_2)$ be open. There exists an equivariant field embedding $E_2\hookrightarrow E_1$ if and only if
\[
 \chi_{E_2/F_2}\circ\alpha=\chi_{E_1/F_1}
 \quad\text{in }\widehat{\Z}^{\times}.
\]
\end{theorem}
We use \cite[Theorem~3.4]{hoshi2025}. Uniqueness of the embedding, when it exists, is \cite[Proposition~2]{uchida1981}.

\subsection{Tame inertia and norms}
Decomposition groups in $P_M$ are the full absolute Galois groups of the corresponding local fields. Here is a justification for using the usual local ramification structure. Fix a prime of $\overline{\Q}$ above a prime of $M$, and form inside a local algebraic closure the union of the completions of finite fields between $M$ and $\Q^s$. This union contains all roots of unity and is closed under radicals: for a fixed exponent, a nonzero local element can be approximated by a global element in the field defining that completion, and their ratio can be made a power. Thus the union has no cyclic extension. Every finite Galois group over a nonarchimedean local field is solvable, by its wild inertia, tame inertia, and unramified filtration. Every closed subgroup of its prosolvable absolute Galois group is prosolvable as well; a nontrivial such group has a nontrivial finite cyclic quotient after passing through a finite solvable quotient. The union is therefore the full local algebraic closure. The usual decomposition-group identification proves the claim.

For an active pair $(v,w)$ from Theorem~\ref{thm:local-correspondence}, and $g\in D_v$, let $b(g)\in\widehat{\Z}$ be its source residue exponent and let $a(g)\in\widehat{\Z}$ be the target residue exponent of $\varphi(g)$. These are measured using $Nv$- and $Nw$-Frobenius, respectively. The target residue exponent is defined on $D_v$ by composition; we do not assume that the map sends all source inertia into target inertia.

\begin{lemma}\label{lem:norm-identity}
For an active pair $(v,w)$ at an odd prime $\ell$,
\begin{equation}\label{eq:norm-identity}
 (Nv)^{b(g)}=(Nw)^{a(g)}\quad\text{in }\Z_\ell^\times
 \qquad(g\in D_v).
\end{equation}
Powers by profinite integers are interpreted by continuous powering in the compact group $\Z_\ell^\times$.
\end{lemma}
\begin{proof}
Let $p\ne\ell$ be the residue characteristic of $v$, and choose a topological generator $\tau$ of a pro-$\ell$ Sylow subgroup of $I_v$. Its image in source tame inertia generates the pro-$\ell$ factor. The tame conjugation relation shows that
\[
 e=g\tau g^{-1}\tau^{-(Nv)^{b(g)}}
\]
belongs to source wild inertia, a pro-$p$ group. In the target tame quotient, the image of $e$ also lies in the pro-$\ell$ factor of tame inertia: both $\varphi(\tau)$ and its conjugate do. Its image is therefore trivial, because it belongs to a pro-$p$ subgroup and to a pro-$\ell$ group. The nontrivial image of $\varphi(\tau)$ in the torsion-free target pro-$\ell$ tame factor then gives \eqref{eq:norm-identity} by the target tame conjugation relation. The target residue characteristic is different from $\ell$ by Theorem~\ref{thm:local-correspondence}.
\end{proof}
